\Askhsh{27092}{4}
\begin{ylh}{1}
    \item 1.2 Συναρτήσεις
    \item 1.8 Συνέχεια συνάρτησης
    \item 2.2 Παραγωγίσιμες συναρτήσεις - Παράγωγος συνάρτηση
    \item 2.3 Κανόνες παραγώγισης
    \item 2.6 Συνέπειες του Θεωρήματος Μέσης Τιμής
    \item 2.7 Τοπικά ακρότατα συνάρτησης.
\end{ylh}
\begin{wrapfigure}[21]{r}{0.45\textwidth}
    \centering
    \begin{tikzpicture}
        \begin{axis}[
            axis lines = middle,
            xlabel = {$x$},
            ylabel = {$y$},
            xmin = -2.5, xmax = 5.5,
            ymin = -3.5, ymax = 2.5,
            xtick = {-2,-1,0,1,2,3,4,5},
            ytick = {-3,-2,-1,0,1,2},
            grid = both,
            grid style = {gray!30},
            samples = 100,
            width = 0.45\textwidth, height = 0.35\textheight,
            tick label style = {font=\scriptsize},
            every axis x label/.style = {at={(current axis.right of origin)}, anchor=north west},
            every axis y label/.style = {at={(current axis.above origin)}, anchor=south east},
        ]
            \addplot[very thick, black, domain=-2.5:5.5] {-x^2 + 3*x};
            \fill (axis cs:0,0) circle (2pt);
            \fill (axis cs:3,0) circle (2pt);
            \node[below right] at (axis cs:4.5,-2.5) {$C_{f'}$};
        \end{axis}
    \end{tikzpicture}
\end{wrapfigure}
Στο παρακάτω σχήμα δίνεται η γραφική παράσταση της παραγώγου $f'$ μιας πολυωνυμικής συνάρτησης $f$ τρίτου βαθμού.
\begin{alist}
    \item Με τη βοήθεια του σχήματος, να μελετήσετε τη συνάρτηση $f$ ως προς τη μονοτονία. \monades{06}
    \item Αν η γραφική παράσταση της $f$ διέρχεται από τα σημεία $A(0,-1)$ και $B(3,2)$, τότε να βρείτε τα ακρότατα της $f$. \monades{04}
    \item Να προσδιορίσετε τον τύπο της $f$. \monades{08}
    \item Να βρείτε το πλήθος ριζών της εξίσωσης $f(x) = \alpha, \alpha \in \mathbb{R}$, στο διάστημα $(0,3)$. \monades{07}
\end{alist}